% =====================================================================
%  Page de garde, remerciements, résumé, tables
% =====================================================================

\begin{titlepage}
\begin{tikzpicture}[remember picture, overlay]
  % Bandeau supérieur
  \fill[medicaDark] (current page.north west) rectangle
    ([yshift=-4.2cm]current page.north east);
  % Liseré d'accent
  \fill[medicaGreen] ([yshift=-4.2cm]current page.north west) rectangle
    ([yshift=-4.35cm]current page.north east);
  % Bandeau inférieur
  \fill[medicaDark] ([yshift=1.5cm]current page.south west) rectangle
    (current page.south east);

  % Marque
  \node[anchor=north west, text=white] at ([shift={(2.6cm,-1.5cm)}]current page.north west)
    {\fontsize{26}{28}\selectfont\bfseries Medicacom};
  \node[anchor=north west, text=medicaSoft] at ([shift={(2.6cm,-2.35cm)}]current page.north west)
    {\small\sffamily AOI TRANSFORMATION};

  % Croix pharmaceutique stylisée
  \begin{scope}[shift={([shift={(-3.4cm,-2.1cm)}]current page.north east)}]
    \fill[medicaGreen, rounded corners=6pt] (-0.9,-0.9) rectangle (0.9,0.9);
    \fill[white] (-0.18,-0.6) rectangle (0.18,0.6);
    \fill[white] (-0.6,-0.18) rectangle (0.6,0.18);
  \end{scope}
\end{tikzpicture}

\vspace*{5.2cm}

\begin{center}
  {\large \varUniversite}\\[3pt]
  {\normalsize \varEcole}\\[3pt]
  {\small \textcolor{medicaMuted}{\varDepartement}}

  \vspace{1.6cm}

  {\small\textcolor{medicaMuted}{\textsc{Rapport de projet}}}\\[2pt]
  {\small\textcolor{medicaMuted}{\varDiplome{} --- \varFormation}}

  \vspace{1.3cm}

  \rule{0.72\textwidth}{0.8pt}

  \vspace{0.9cm}

  {\LARGE\bfseries\color{medicaDeep} \varTitre\par}

  \vspace{0.7cm}

  {\large\itshape\color{medicaMuted} \varSousTitre\par}

  \vspace{0.9cm}

  \rule{0.72\textwidth}{0.8pt}

  \vspace{2.1cm}

  \begin{minipage}{0.44\textwidth}
    \raggedright
    \textbf{Réalisé par}\\[4pt]
    \varAuteur
  \end{minipage}
  \hfill
  \begin{minipage}{0.44\textwidth}
    \raggedleft
    \textbf{Encadré par}\\[4pt]
    \varEncadrantAcademique\\
    \varEncadrantEntreprise
  \end{minipage}

  \vspace{1.5cm}

  \textbf{Entreprise d'accueil :} \varEntreprise

  \vfill

  {\color{white}\textbf{Année universitaire \varAnnee}}
  \vspace{0.35cm}
\end{center}
\end{titlepage}

% ---------------------------------------------------------------- remerciements
\chapter*{Remerciements}
\addcontentsline{toc}{chapter}{Remerciements}
\markboth{Remerciements}{}

Je tiens à remercier \varEncadrantAcademique, dont l'encadrement académique a
structuré la conduite de ce travail, ainsi que \varEncadrantEntreprise, pour
l'accompagnement quotidien au sein de \varEntreprise{} et la confiance accordée sur un
sujet aussi transverse que le pilotage de l'adoption de l'intelligence artificielle.

Mes remerciements vont également aux équipes métier qui ont accepté de décrire leurs
processus, leurs irritants et leurs indicateurs. Sans ces échanges, le référentiel de
diagnostic et le questionnaire de qualification n'auraient été qu'un exercice
théorique.

Je remercie enfin les membres du jury pour le temps consacré à l'évaluation de ce
travail.

% ---------------------------------------------------------------------- résumé
\chapter*{Résumé}
\addcontentsline{toc}{chapter}{Résumé}
\markboth{Résumé}{}

Les laboratoires pharmaceutiques accumulent des idées d'intelligence artificielle sans
disposer d'un cadre commun pour décider lesquelles financer. Les besoins remontent de
départements cloisonnés, les données ne sont pas toujours exploitables, la conformité
réglementaire pèse sur chaque décision, et la valeur créée reste rarement mesurée après
la mise en production.

Ce projet répond à ce problème par une plateforme SaaS multi-entreprises,
\appname, qui couvre l'ensemble du parcours : diagnostic de maturité inspiré de
la démarche TOGAF, qualification de la solution la plus adaptée, évaluation du
readiness, priorisation du portefeuille, génération de la fiche projet, suivi du cycle
de vie et mesure du retour sur investissement.

La plateforme repose sur une API REST écrite en Python avec FastAPI et SQLAlchemy~2,
une base PostgreSQL et une interface React en TypeScript. Le contrôle d'accès associe
41~permissions granulaires, 11~rôles fonctionnels et cinq niveaux de périmètre, avec
une isolation stricte entre organisations clientes. Chaque décision produite par le
système est explicable : le score de priorité renvoie la contribution de chacun de ses
six critères, et la recommandation de solution expose les règles satisfaites, les
alternatives écartées et les garde-fous à mettre en place.

Un point mérite d'être énoncé sans détour : la plateforme ne contient aujourd'hui
\emph{aucun modèle d'apprentissage automatique}. Son intelligence est celle d'un moteur
de règles pondérées, déterministe et auditable. Ce choix est assumé et documenté au
chapitre~\ref{chap:realisation}; l'intégration d'un modèle de langage figure parmi les
perspectives, non parmi les fonctionnalités livrées.

\vspace{0.8em}
\noindent\textbf{Mots-clés} — adoption de l'IA, industrie pharmaceutique, TOGAF,
readiness, portefeuille de projets, moteur de règles, gouvernance, multi-tenant,
FastAPI, React, PostgreSQL.

% -------------------------------------------------------------------- abstract
\chapter*{Abstract}
\addcontentsline{toc}{chapter}{Abstract}
\markboth{Abstract}{}

Pharmaceutical companies accumulate artificial intelligence ideas without a shared
framework for deciding which ones deserve funding. Needs come from siloed departments,
data is not always usable, regulatory compliance weighs on every decision, and the
value actually delivered is rarely measured once a system reaches production.

This project addresses that problem with \appname, a multi-tenant SaaS platform
covering the full journey: a TOGAF-inspired maturity assessment, qualification of the
most suitable solution type, readiness scoring, portfolio prioritisation, project
charter generation, lifecycle tracking and return-on-investment measurement.

The platform is built on a REST API written in Python with FastAPI and SQLAlchemy~2, a
PostgreSQL database and a React frontend in TypeScript. Access control combines
41~granular permissions, 11~functional roles and five scope levels, with strict
isolation between client organisations. Every decision the system produces is
explainable: the priority score returns the contribution of each of its six criteria,
and the solution recommendation exposes the rules that fired, the alternatives that
were ruled out and the guardrails to put in place.

One point deserves to be stated plainly: the platform currently contains
\emph{no machine learning model}. Its intelligence is that of a weighted, deterministic
and auditable rules engine. This is a deliberate choice, documented in
Chapter~\ref{chap:realisation}; integrating a large language model belongs to the
outlook, not to the delivered features.

\vspace{0.8em}
\noindent\textbf{Keywords} — AI adoption, pharmaceutical industry, TOGAF, readiness,
project portfolio, rules engine, governance, multi-tenancy, FastAPI, React, PostgreSQL.

% ----------------------------------------------------------------------- tables
\cleardoublepage
\renewcommand{\contentsname}{Table des matières}
\tableofcontents

\cleardoublepage
\listoffigures
\addcontentsline{toc}{chapter}{Table des figures}

\cleardoublepage
\listoftables
\addcontentsline{toc}{chapter}{Liste des tableaux}
